\sectionheading{5.6 উচ্চস্তরের ভাষা (High Level Language)}

Machine Language ও Assembly Language কম্পিউটারের Hardware-এর খুব
কাছাকাছি হলেও মানুষের জন্য এসব ভাষায় Program লেখা তুলনামূলক কঠিন।
বিশেষ করে Processor, Register, Memory Address এবং Machine-specific
Instruction সম্পর্কে বিস্তারিত ধারণা প্রয়োজন হয়।

এই সমস্যা কমিয়ে মানুষের জন্য সহজ, বোধগম্য ও দ্রুত Program লেখার
উদ্দেশ্যে এমন Programming Language তৈরি করা হয়, যার শব্দ ও গঠন
মানুষের পরিচিত ভাষা ও গাণিতিক প্রকাশের তুলনামূলক কাছাকাছি।
এ ধরনের ভাষাকে \textbf{High Level Language} বলা হয়।

%================================================
% Definition
%================================================
\booksubheading{উচ্চস্তরের ভাষা কাকে বলে?}

\begin{definitionbox}

\textbf{\color{darkTeal}সংজ্ঞা}

যে Programming Language-এর \textbf{শব্দ, প্রতীক ও নির্দেশ মানুষের
ভাষা ও গাণিতিক প্রকাশের তুলনামূলক কাছাকাছি}, ফলে সহজে Program
লেখা, পড়া ও বোঝা যায়, তাকে
\textbf{উচ্চস্তরের ভাষা (High Level Language)} বলে।

\end{definitionbox}

\textbf{সহজভাবে মনে রাখবে—}

\begin{center}

\textbf{
মানুষের জন্য সহজ ও বোধগম্য Programming Language
\(\;=\;\)
High Level Language
}

\end{center}

High Level Language ব্যবহার করার সময় Programmer-কে সাধারণত
Processor-এর Binary Instruction বা Hardware-এর অভ্যন্তরীণ গঠন
নিয়ে সরাসরি কাজ করতে হয় না।

%================================================
% Why easy
%================================================
\booksubheading{High Level Language কেন মানুষের জন্য সহজ?}

High Level Language-এ সাধারণত অর্থবোধক Keyword, Variable,
Mathematical Operator, Conditional Statement, Loop, Function
ইত্যাদি ব্যবহার করা হয়।

\begin{examplebox}{সহজ উদাহরণ}

ধরা যাক—

\begin{center}

{\textenglish{\ttfamily total = a + b}}

\end{center}

Instructionটি দেখেই সহজে বোঝা যায় যে
\textenglish{\ttfamily a} ও \textenglish{\ttfamily b}-এর মান যোগ করে
ফলাফল \textenglish{\ttfamily total}-এ রাখা হচ্ছে।

Machine Language-এ একই ধরনের কাজ 0 ও 1-এর Binary Instruction
দিয়ে প্রকাশ করতে হয়।

\end{examplebox}

%================================================
% Examples
%================================================
\booksubheading{High Level Language-এর উদাহরণ}

High Level Language-এর পরিচিত উদাহরণ হলো—

\begin{center}

\textbf{
Python, Java, C++, BASIC, Pascal, FORTRAN
}

\end{center}

নতুন সিলেবাস অনুযায়ী \textbf{Java, Oracle ও Python} সম্পর্কে
পরবর্তী অংশে আলাদাভাবে আলোচনা করা হবে।

%================================================
% Features
%================================================
\booksubheading{High Level Language-এর বৈশিষ্ট্য}

\begin{tabularx}{\linewidth}{|>{\bfseries}p{0.27\linewidth}|Y|}
\hline

\rowcolor{softBlue}
\multicolumn{1}{|c|}{\textbf{বৈশিষ্ট্য}}
&
\multicolumn{1}{c|}{\textbf{ব্যাখ্যা}}
\\
\hline

সহজ ও বোধগম্য
&
Program লেখা, পড়া ও বোঝা তুলনামূলক সহজ।
\\
\hline

Problem-oriented
&
Hardware-এর সূক্ষ্ম বিষয় অপেক্ষা Problem Solving-এর ওপর বেশি
গুরুত্ব দেওয়া যায়।
\\
\hline

Machine Independent
&
নির্দিষ্ট Hardware-এর ওপর সরাসরি নির্ভরশীলতা তুলনামূলক কম।
\\
\hline

Translator প্রয়োজন
&
কম্পিউটার High Level Code সরাসরি বুঝতে পারে না।
\\
\hline

সহজ Syntax
&
অর্থবোধক Keyword ও পরিচিত Mathematical Expression ব্যবহার করা যায়।
\\
\hline

Debugging সহজ
&
Error খুঁজে বের ও সংশোধন তুলনামূলক সহজ।
\\
\hline

Portability বেশি
&
প্রয়োজনীয় Translator/Runtime থাকলে একই Source Program ভিন্ন
System-এ ব্যবহার করা তুলনামূলক সহজ।
\\
\hline

Library সুবিধা
&
অনেক High Level Language-এ Built-in Function ও Library পাওয়া যায়।
\\
\hline

\end{tabularx}

%================================================
% Translation
%================================================
\booksubheading{High Level Language কীভাবে কম্পিউটারে চলে?}

কম্পিউটার High Level Language সরাসরি বুঝতে পারে না।
তাই High Level Language-এ লেখা Program-কে কম্পিউটারের বোধগম্য
রূপে পরিবর্তন করার জন্য \textbf{Translator Program} প্রয়োজন হয়।

\begin{center}
\begin{adjustbox}{max width=0.98\linewidth}
\fcolorbox{darkBlue}{softBlue}{
    \parbox[c][8mm][c]{0.25\linewidth}{
        \centering\textbf{High Level Program}
    }
}
\quad
\(\longrightarrow\)
\quad
\fcolorbox{mainTeal}{softGreen}{
    \parbox[c][8mm][c]{0.20\linewidth}{
        \centering\textbf{Compiler / Interpreter}
    }
}
\quad
\(\longrightarrow\)
\quad
\fcolorbox{darkBlue}{softBlue}{
    \parbox[c][8mm][c]{0.27\linewidth}{
        \centering\textbf{Machine-understandable Code}
    }
}
\end{adjustbox}
\end{center}

Compiler ও Interpreter সম্পর্কে Translator Program অংশে
বিস্তারিত আলোচনা করা হবে।

%================================================
% Source Program
%================================================
\booksubheading{Source Program কী?}

High Level Language-এ Programmer যে মূল Program লেখে,
তাকে সাধারণভাবে \textbf{Source Program বা Source Code} বলা হয়।

\begin{definitionbox}

\textbf{\color{darkTeal}Source Program}

Programmer কোনো Programming Language ব্যবহার করে যে মূল
Program বা Code লেখে, তাকে \textbf{Source Program} বলে।

\end{definitionbox}

\begin{center}

\textbf{
Programmer-এর লেখা Code
\(\;=\;\)
Source Program
}

\end{center}

Translator এই Source Program-কে Machine-এর উপযোগী রূপে
রূপান্তর করে।

\begin{rememberbox}

\textbf{\color{deepOrange}মনে রাখবে}

High Level Language-এ লেখা মূল Program হলো
\textbf{Source Program বা Source Code}।

অনুবাদের পর Machine-এর উপযোগী Code-কে প্রসঙ্গ অনুযায়ী
Target Code বা Object Code বলা হতে পারে।

\end{rememberbox}

%================================================
% Advantages
%================================================
\booksubheading{High Level Language-এর সুবিধা}

\begin{compacttwocolpoints}

\item Program লেখা তুলনামূলক \textbf{সহজ ও দ্রুত}।

\item Program পড়া ও Logic বুঝতে সহজ।

\item Machine বা Assembly Language-এর তুলনায় কম Instruction-এ
বড় কাজ করা যায়।

\item Program-এ Error হওয়ার সম্ভাবনা তুলনামূলক কম।

\item Error খুঁজে বের ও সংশোধন বা \textbf{Debugging} সহজ।

\item Program পরিবর্তন ও Maintenance সহজ।

\item Hardware-এর সূক্ষ্ম গঠন সম্পর্কে গভীর ধারণা ছাড়াই
সাধারণ Program লেখা যায়।

\item Built-in Function ও Library ব্যবহার করে দ্রুত Program তৈরি করা যায়।

\item Program ভিন্ন Machine/System-এ ব্যবহার করা তুলনামূলক সহজ।

\item বড় ও জটিল Application তৈরি করা সুবিধাজনক।

\end{compacttwocolpoints}

%================================================
% Disadvantages
%================================================
\booksubheading{High Level Language-এর অসুবিধা}

\begin{itemize}

\item কম্পিউটার সরাসরি High Level Language বুঝতে পারে না।

\item Machine Code-এ রূপান্তরের জন্য
\textbf{Compiler বা Interpreter} প্রয়োজন হয়।

\item Low Level Language-এর তুলনায় Hardware-এর ওপর সরাসরি
নিয়ন্ত্রণ কম।

\item অনেক ক্ষেত্রে Low Level Program-এর তুলনায় বেশি Memory বা
System Resource প্রয়োজন হতে পারে।

\item Hardware-specific কাজের ক্ষেত্রে Low Level Language-এর মতো
সরাসরি নিয়ন্ত্রণ নাও পাওয়া যেতে পারে।

\item Program লেখার জন্য সংশ্লিষ্ট Language-এর Syntax ও নিয়ম জানতে হয়।

\end{itemize}

%================================================
% Uses
%================================================
\booksubheading{High Level Language-এর ব্যবহার}

High Level Language বিভিন্ন ধরনের Software Development ও
Problem Solving-এ ব্যবহার করা হয়। যেমন—

\begin{compacttwocolpoints}

\item Application Program তৈরিতে।

\item Data Processing-এর কাজে।

\item Web ও Software Development-এ।

\item বৈজ্ঞানিক ও গাণিতিক Problem সমাধানে।

\item Business Application তৈরিতে।

\item Database-ভিত্তিক System তৈরিতে।

\item Automation-এর কাজে।

\item শিক্ষা ও গবেষণায়।

\item Data Analysis ও Artificial Intelligence-এর মতো আধুনিক ক্ষেত্রে।

\end{compacttwocolpoints}

%================================================
% Comparison
%================================================
\booksubheading{Machine Language ও High Level Language-এর পার্থক্য}

\begin{tabularx}{\linewidth}{|>{\bfseries}p{0.20\linewidth}|Y|Y|}
\hline

\rowcolor{softBlue}
\multicolumn{1}{|c|}{\textbf{বিষয়}}
&
\multicolumn{1}{c|}{\textbf{Machine Language}}
&
\multicolumn{1}{c|}{\textbf{High Level Language}}
\\
\hline

ভাষার রূপ
&
0 ও 1-এর Binary Code
&
অর্থবোধক শব্দ, প্রতীক ও Statement
\\
\hline

মানুষের জন্য
&
লেখা ও বোঝা কঠিন
&
তুলনামূলক সহজ
\\
\hline

কম্পিউটার
&
সরাসরি বুঝতে পারে
&
সরাসরি বুঝতে পারে না
\\
\hline

Translator
&
প্রয়োজন নেই
&
Compiler বা Interpreter প্রয়োজন
\\
\hline

Machine Dependence
&
বেশি
&
তুলনামূলক কম
\\
\hline

Program লেখা
&
কঠিন ও সময়সাপেক্ষ
&
সহজ ও দ্রুত
\\
\hline

Error শনাক্ত
&
কঠিন
&
তুলনামূলক সহজ
\\
\hline

Debugging
&
কঠিন
&
সহজ
\\
\hline

Hardware Knowledge
&
বেশি প্রয়োজন
&
তুলনামূলক কম প্রয়োজন
\\
\hline

Portability
&
কম
&
তুলনামূলক বেশি
\\
\hline

\end{tabularx}

%================================================
% Machine Independence
%================================================
\booksubheading{High Level Language কেন Machine Independent বলা হয়?}

High Level Language-এ Program লেখার সময় Programmer কোনো একটি
নির্দিষ্ট Processor-এর Binary Instruction ব্যবহার করে না।

তাই একই Source Program প্রয়োজনীয় Translator বা Runtime-এর সাহায্যে
ভিন্ন System-এ ব্যবহার করা তুলনামূলক সহজ।

\begin{examplebox}{সহজভাবে বোঝো}

একটি Python Source Program Programmer একইভাবে লিখতে পারে।
ভিন্ন Operating System বা Machine-এ প্রয়োজনীয় Python Environment
থাকলে সেই Program ব্যবহার করা সম্ভব হতে পারে।

অর্থাৎ Source Code সরাসরি কোনো একটি নির্দিষ্ট Processor-এর
Binary Instruction-এর ওপর নির্ভর করে না।

\end{examplebox}

\begin{rememberbox}

\textbf{\color{deepOrange}সতর্কতা}

Machine Independent বলার অর্থ এই নয় যে একই Executable File
সব Machine-এ সরাসরি চলবে।

প্রয়োজনীয় \textbf{Translator, Runtime বা Platform Support}
থাকতে হয়।

\end{rememberbox}

%================================================
% Why Popular
%================================================
\booksubheading{High Level Language কেন জনপ্রিয়?}

High Level Language Programmer-কে Hardware-এর জটিল Binary
Instruction বা Register নিয়ে সরাসরি কাজ করার প্রয়োজন অনেকটাই কমিয়ে দেয়।

ফলে Programmer বেশি গুরুত্ব দিতে পারে—

\begin{compacttwocolpoints}

\item Problem বোঝায়,

\item Logic তৈরিতে,

\item Algorithm তৈরি করতে,

\item Program দ্রুত Development করতে,

\item Testing ও Debugging করতে।

\end{compacttwocolpoints}

এ কারণেই আধুনিক Software Development-এ High Level Language
ব্যাপকভাবে ব্যবহৃত হয়।

%================================================
% Generation
%================================================
\booksubheading{High Level Language ও Generation}

Programming Language-এর Generation-ভিত্তিক আলোচনায়
\textbf{Third Generation Language (3GL)}-এর অধিকাংশ ভাষাই
High Level Language।

যেমন—

\begin{center}

\textbf{
FORTRAN, BASIC, Pascal, C
}

\end{center}

3GL-এর পরিচিত উদাহরণ হিসেবে প্রচলিত পাঠ্যবইয়ে আলোচনা করা হয়।

তবে—

\begin{center}

\textbf{
High Level Language
\(\;\neq\;\)
Third Generation Language (3GL)
}

\end{center}

দুটি সম্পূর্ণ একই শ্রেণিবিভাগ নয়।

High Level Language মূলত ভাষার \textbf{সহজবোধ্যতা ও Hardware-এর
সঙ্গে সম্পর্ক} বোঝায়, আর 3GL হলো Programming Language-এর
\textbf{Generation-ভিত্তিক শ্রেণিবিভাগ}।

%================================================
% Final Reminder
%================================================
\begin{rememberbox}

\textbf{\color{deepOrange}মনে রাখবে}

\begin{itemize}

\item High Level Language মানুষের জন্য
\textbf{সহজ ও বোধগম্য}।

\item এটি Hardware-এর পরিবর্তে
\textbf{Problem Solving-এর কাছাকাছি}।

\item High Level Program কম্পিউটার
\textbf{সরাসরি বুঝতে পারে না}।

\item তাই \textbf{Compiler বা Interpreter} প্রয়োজন হয়।

\item Programmer-এর লেখা মূল Code হলো
\textbf{Source Program বা Source Code}।

\item High Level Language সাধারণভাবে
\textbf{Machine Independent} এবং Portability বেশি।

\item Program লেখা, Error শনাক্ত ও Debugging তুলনামূলক সহজ।

\item Python ও Java হলো High Level Language-এর পরিচিত উদাহরণ।

\item High Level Language এবং 3GL
\textbf{সম্পূর্ণ একই শ্রেণিবিভাগ নয়}।

\end{itemize}

\end{rememberbox}

\small

%================================================
% 5.7 Java
%================================================
